\subsection{The valuation-law gate: Tao's Proposition~1.9}

The proposition numbered 1.9 in the controlling v7 text is the distribution
theorem for the finite Syracuse valuation tuple, not the later first-passage
stabilisation theorem.  Its statement is at v7 source lines~269--278 and its
proof occupies lines~597--647 \cite{Tao2026v7}.  Proposition~1.11 is the
first-passage result.  This numbering matters in the dependency graph:
Proposition~1.14 and Proposition~1.9 are the two inputs used to obtain
Proposition~1.11.

The frozen routes identify the v5 and journal manifestations of the same
work.  Direct comparison uses v5 source lines~269--278 and~591--641, v7
source lines~269--278 and~597--647, and journal physical pages~6 and~17--18.
Routing metadata is not used as mathematical evidence.

\begin{sourceissue}[three-manifestation defects in Proposition~1.9]
\label{sourceissue:Tao-Prop19-three-manifestations}
All three inspected manifestations print the same substantive residue-fibre
argument and preserve three operative presentation defects.  In the
statement, the ambient ring is printed as \(\mathbb Z/2^\ell\mathbb Z\),
although \(\ell\) is undefined; both the displayed law and every subsequent
use require \(\mathbb Z/2^{n'}\mathbb Z\).  In the tail proof, one intermediate
line changes the first-crossing event
\(A_k<n'\leq A_{k+1}\) into \(A_k\leq n'<A_{k+1}\); the former is the event
used before and after that line.  Finally, the two central truncations are
printed as \(|\vec a|<m\), although no \(m\) is defined in that section; the
proof requires \(|\vec a|<n'\).

The source also invokes the empty prefix sum \(A_0\) without extending its
earlier nonempty-sum convention, and it does not dispatch \(n=0\).  In the
v5/v7 TeX proof of the prerequisite valuation-uniqueness lemma, a function is
once written without its argument and \(\operatorname{Syr}^{N-1}(N)\) is
printed where the induction requires \(\operatorname{Syr}^{n-1}(N)\).
The reconstruction below supplies all of these data explicitly.
\end{sourceissue}

For \(s\geq1\), put
\[
 \mathcal O_s=(2\mathbb Z+1)/2^s\mathbb Z.
\]
Also set \(\mathcal O_0=\mathbb Z/\mathbb Z\), the singleton residue space,
so that the \(n=0,n'=0\) endpoint in the source statement is typed.  For
\(s\geq1\), \(\mathcal O_s\) has \(2^{s-1}\) elements.  For
\(1\leq t\leq s\), let
\[
 \pi_{s,t}:\mathcal O_s\longrightarrow\mathcal O_t,
 \qquad [u]_{2^s}\longmapsto[u]_{2^t}
\]
be reduction.  Every fibre has exactly \(2^{s-t}\) elements, so
\((\pi_{s,t})_*\operatorname{Unif}(\mathcal O_s)\) is
\(\operatorname{Unif}(\mathcal O_t)\).  For signed measures \(\mu,\nu\) on
\(\mathcal O_s\), the source's unhalved total-variation convention gives the
exact contraction
\begin{equation}
\label{eq:Tao-Prop19-TV-contraction}
\begin{split}
 \| (\pi_{s,t})_*\mu-(\pi_{s,t})_*\nu\|_1
 & =\sum_{x\in\mathcal O_t}
 \left|\sum_{u\in\pi_{s,t}^{-1}(x)}(\mu(u)-\nu(u))\right|\\
 &\leq\sum_{u\in\mathcal O_s}|\mu(u)-\nu(u)|.
\end{split}
\end{equation}

\begin{lemma}[valuation prefix and its exact residue fibre]
\label{lem:Tao-Prop19-residue-fibre}
Let \(n\geq1\), let \(N\in2\mathbb N+1\), and let
\(\vec a=(a_1,\ldots,a_n)\in(\mathbb N+1)^n\).  Define
\[
 A_0=0,\qquad A_j=a_1+\cdots+a_j\quad(1\leq j\leq n),
 \qquad A=A_n,
\]
and
\[
 C_n(\vec a)=\sum_{i=1}^n3^{n-i}2^{A_{i-1}}.
\]
Then
\begin{equation}
\label{eq:Tao-Prop19-affine-numerator}
 \operatorname{Aff}_{\vec a}(N)
 =\frac{3^nN+C_n(\vec a)}{2^A},
\end{equation}
and
\begin{equation}
\label{eq:Tao-Prop19-prefix-equivalence}
 \vec a^{(n)}(N)=\vec a
 \quad\Longleftrightarrow\quad
 \operatorname{Aff}_{\vec a}(N)\in2\mathbb N+1.
\end{equation}
Consequently the event in \eqref{eq:Tao-Prop19-prefix-equivalence} is the
single odd residue cylinder
\begin{equation}
\label{eq:Tao-Prop19-cylinder}
 N\equiv b_{\vec a}\pmod{2^{A+1}},
 \qquad
 b_{\vec a}\equiv
 3^{-n}\bigl(2^A-C_n(\vec a)\bigr)\pmod{2^{A+1}}.
\end{equation}
If \(A<n'\), the inverse image of this cylinder in \(\mathcal O_{n'}\) has
exactly \(2^{n'-A-1}\) elements and hence uniform mass \(2^{-A}\).
\end{lemma}

\begin{proof}
Multiplying the source's affine-offset formula by \(2^A\) gives
\eqref{eq:Tao-Prop19-affine-numerator}.  The forward implication in
\eqref{eq:Tao-Prop19-prefix-equivalence} is the defining Syracuse iteration.
For the converse, induct on \(n\).  When \(n=1\), the equality
\(
  2^{a_1}\operatorname{Aff}_{a_1}(N)=3N+1
\)
and the oddness of \(\operatorname{Aff}_{a_1}(N)\) give
\(a_1=\nu_2(3N+1)\) directly.  Now let \(n\geq2\) and write
\[
 X=\operatorname{Aff}_{a_1,\ldots,a_{n-1}}(N),
 \qquad
 Y=\operatorname{Aff}_{\vec a}(N)\in2\mathbb N+1.
\]
Both affine maps preserve positivity and \(X\in\mathbb Z[1/2]\).  The last
step gives
\[
 2^{a_n}Y=3X+1.
\]
Thus \(3X\) is an odd integer.  A rational number in
\(\mathbb Z[1/2]\) whose product by \(3\) is integral is itself integral;
hence \(X\) is a positive odd integer.  The induction hypothesis identifies
the first \(n-1\) valuation coordinates.  Since
\((3X+1)/2^{a_n}=Y\) is odd, \(a_n=\nu_2(3X+1)\), proving the last coordinate.

An integer quotient in \eqref{eq:Tao-Prop19-affine-numerator} is odd exactly
when its numerator is congruent to \(2^A\) modulo \(2^{A+1}\).  Because
\(3^n\) is a unit modulo \(2^{A+1}\), this gives the unique residue in
\eqref{eq:Tao-Prop19-cylinder}.  It is odd: the first term of
\(C_n(\vec a)\) is \(3^{n-1}\), every other term is even, and \(2^A\) is
even.  Finally, reduction
\(\mathcal O_{n'}\to\mathcal O_{A+1}\) has fibres of size
\(2^{n'-(A+1)}\), proving the last assertion.
\end{proof}

The exponential estimate used in both tails and in the accumulated cylinder
error can be written without an unnamed appeal to Stirling's formula.  Let
\[
 H(p)=-p\log p-(1-p)\log(1-p),
 \qquad
 p_0=\frac1{2+c_0/2}<\frac12,
 \qquad
 \eta(c_0)=\log2-H(p_0)>0.
\]
If \(M=n'-1\), \(n'\geq(2+c_0)n\), and \(n\geq2/c_0\), then
\(n/M\leq p_0\).  Set \(z_0=p_0/(1-p_0)\in(0,1)\).  Since
\(z_0^k\geq z_0^n\) for \(0\leq k\leq n\),
\begin{equation}
\label{eq:Tao-Prop19-binomial-bound}
 2^{-M}\sum_{k=0}^n\binom Mk
 \leq z_0^{-n}\left(\frac{1+z_0}{2}\right)^M
 \leq \exp\bigl(-\eta(c_0)M\bigr).
\end{equation}
The finitely many positive \(n<2/c_0\) are absorbed into a constant depending
on \(c_0\).  This proves a bound \(C_{c_0}2^{-c_1n}\), for example with any
fixed \(0<c_1<\eta(c_0)/\log2\).

\begin{proposition}[Tao Proposition~1.9, residue-fibre reconstruction]
\label{prop:Tao-v7-Prop19-valuation-law}
Let \(n\in\mathbb N\), let \(\mathbf N\) take values in
\(2\mathbb N+1\), let \(c_0>0\), and let
\(n'\in\mathbb N\) satisfy \(n'\geq(2+c_0)n\).  Suppose, with the unhalved
total variation used by the source, that
\begin{equation}
\label{eq:Tao-Prop19-hypothesis}
 d_{\mathrm{TV}}\left(
 \mathbf N\bmod2^{n'},\operatorname{Unif}(\mathcal O_{n'})
 \right)\leq K2^{-n'}
\end{equation}
for a hypothesis constant \(K=K(c_0)\geq0\), independent of
\(n,n'\), and the law of \(\mathbf N\).  Then there are
\(c_1=c_1(c_0)>0\) and \(C=C(c_0)>0\) such that
\begin{equation}
\label{eq:Tao-Prop19-conclusion}
 d_{\mathrm{TV}}\left(
 \vec a^{(n)}(\mathbf N),\operatorname{Geom}(2)^n
 \right)\leq C(1+K)2^{-c_1n}.
\end{equation}
Thus the source statement follows with exactly the source-permitted
\(c_0\)-dependence of the implicit constants.
\end{proposition}

\begin{proof}
For \(n=0\), both laws are the point mass at the empty tuple, so the distance
is zero.  Assume \(n\geq1\), and write
\[
 \vec a^{(n)}(\mathbf N)=(\mathsf a_1,\ldots,\mathsf a_n),
 \qquad
 S_0=0,\qquad S_j=\mathsf a_1+\cdots+\mathsf a_j.
\]
Because every \(\mathsf a_j\) is positive, the actual-valuation tail is the
disjoint first-crossing union
\begin{equation}
\label{eq:Tao-Prop19-first-crossing}
 \{S_n\geq n'\}
 =\bigsqcup_{k=0}^{n-1}\{S_k<n'\leq S_{k+1}\}.
\end{equation}
 Fix \(k\) and a positive prefix \((a_1,\ldots,a_k)\), put
 \(A_0=0\) and \(A_j=\sum_{r=1}^{j}a_r\) for \(1\leq j\leq k\), and assume
 \(A_k<n'\).
On the corresponding crossing event, the exact affine numerator at time
\(k+1\) is divisible by \(2^{S_{k+1}}\), hence
\begin{equation}
\label{eq:Tao-Prop19-crossing-residue}
 3^{k+1}\mathbf N+
 \sum_{i=1}^{k+1}3^{k+1-i}2^{A_{i-1}}
 \equiv0\pmod{2^{n'}}.
\end{equation}
The displayed sum depends only on the fixed \(k\)-prefix, including
\(A_0=0\).  Its first term is odd and all later terms are even, so
\eqref{eq:Tao-Prop19-crossing-residue} fixes one odd residue of
\(\mathbf N\) modulo \(2^{n'}\).  Its probability is at most
\((2+K)2^{-n'}\): the uniform mass of one element of \(\mathcal O_{n'}\) is
\(2^{1-n'}\), and \eqref{eq:Tao-Prop19-hypothesis} controls the error.

Positive \(k\)-tuples with \(A_k<n'\) correspond bijectively to the
\(k\)-element subsets
\[
 \{A_1<\cdots<A_k\}\subseteq\{1,\ldots,n'-1\},
\]
and therefore number \(\binom{n'-1}{k}\), including the one empty prefix at
\(k=0\).  Summing \eqref{eq:Tao-Prop19-first-crossing} and applying
\eqref{eq:Tao-Prop19-binomial-bound} gives
\begin{equation}
\label{eq:Tao-Prop19-actual-tail}
 \mathbb P(S_n\geq n')\ll_{c_0}(1+K)2^{-c_1n}.
\end{equation}

For the model law, realize independent \(\operatorname{Geom}(2)\)
coordinates as waiting times between successive successes in fair Bernoulli
trials.  Their sum is at least \(n'\) exactly when the first \(n'-1\) trials
contain at most \(n-1\) successes.  Hence
\begin{equation}
\label{eq:Tao-Prop19-model-tail}
 \mathbb P\left(|\operatorname{Geom}(2)^n|\geq n'\right)
 =2^{-(n'-1)}\sum_{k=0}^{n-1}\binom{n'-1}{k}
 \ll_{c_0}2^{-c_1n}.
\end{equation}

It remains to compare the two laws on positive tuples
\(\vec a=(a_1,\ldots,a_n)\) with \(A=|\vec a|<n'\).  Lemma
\ref{lem:Tao-Prop19-residue-fibre}, the exact pushforward
\eqref{eq:Tao-Prop19-TV-contraction}, and
\eqref{eq:Tao-Prop19-hypothesis} give
\begin{equation}
\label{eq:Tao-Prop19-point-mass}
 \left|
 \mathbb P(\vec a^{(n)}(\mathbf N)=\vec a)-2^{-A}
 \right|\leq K2^{-n'}.
\end{equation}
The term \(2^{-A}\) is exactly the mass of \(\vec a\) under
\(\operatorname{Geom}(2)^n\).  Positive \(n\)-tuples with \(A<n'\) are in
bijection with \(n\)-element subsets of \(\{1,\ldots,n'-1\}\), so their
number is \(\binom{n'-1}{n}\).  Thus their accumulated error is bounded by
\[
 K2^{-n'}\binom{n'-1}{n}\ll_{c_0}K2^{-c_1n}
\]
using \eqref{eq:Tao-Prop19-binomial-bound}.  Adding the two tail bounds
\eqref{eq:Tao-Prop19-actual-tail} and
\eqref{eq:Tao-Prop19-model-tail} is justified by the exact decomposition
\[
\begin{split}
 d_{\mathrm{TV}}\!\left(
 \vec a^{(n)}(\mathbf N),\operatorname{Geom}(2)^n\right)
 &\leq
 \sum_{\substack{\vec a\in(\mathbb N+1)^n\\|\vec a|<n'}}
 \left|\mathbb P(\vec a^{(n)}(\mathbf N)=\vec a)-2^{-|\vec a|}\right|\\
 &\quad+
 \mathbb P(S_n\geq n')
 +\mathbb P\!\left(|\operatorname{Geom}(2)^n|\geq n'\right).
\end{split}
\]
This proves \eqref{eq:Tao-Prop19-conclusion}.
\end{proof}

The proof exposes the exact morphism suppressed in the printed argument.
Define
\[
 \mathcal U_{n,n'}=
 \left\{u\in\mathcal O_{n'}:
 \begin{array}{l}
 \text{some positive odd representative \(N\) of \(u\) satisfies}\\[-2pt]
 |\vec a^{(n)}(N)|<n'
 \end{array}\right\}.
\]
Lemma~\ref{lem:Tao-Prop19-residue-fibre} shows that the displayed condition
and the resulting tuple are independent of the chosen positive
representative: a tuple of total \(A<n'\) is determined already modulo
\(2^{A+1}\).  Hence
\[
 V_{n,n'}:\mathcal U_{n,n'}\longrightarrow\mathcal D_{n,n'},
 \qquad
 V_{n,n'}([N]_{2^{n'}})=\vec a^{(n)}(N),
\]
is a well-defined map, where
\[
 \mathcal D_{n,n'}=
 \{\vec a\in(\mathbb N+1)^n:|\vec a|<n'\}.
\]
It is surjective, and its fibre over \(\vec a\) is precisely the cylinder
\(b_{\vec a}\pmod{2^{A+1}}\), containing
\(2^{n'-A-1}\) odd residues modulo \(2^{n'}\).  Its set-theoretic kernel pair
is
\[
 \mathcal U_{n,n'}\mathbin{\times_{\mathcal D_{n,n'}}}\mathcal U_{n,n'}
 =\{(u,v):V_{n,n'}(u)=V_{n,n'}(v)\},
\]
namely equality of the first \(n\) valuation coordinates.  There is an
explicit section: for each \(\vec a\), take the least positive odd
representative \(\widetilde b_{\vec a}\) of the cylinder
\(b_{\vec a}\pmod{2^{A+1}}\) and put
\[
 s_{n,n'}(\vec a)=[\widetilde b_{\vec a}]_{2^{n'}},
 \qquad V_{n,n'}\circ s_{n,n'}=\operatorname{id}_{\mathcal D_{n,n'}}.
\]
This section chooses all unrecovered higher binary digits to be zero.  The map
has no left inverse, and therefore no two-sided inverse: the tuple
\((1,\ldots,1)\) has \(A=n\) and fibre size
\(2^{n'-n-1}\geq2\), because
\(n'\geq(2+c_0)n\) with \(n\geq1\).

\begin{nonclaim}
Proposition~\ref{prop:Tao-v7-Prop19-valuation-law} certifies only the
valuation-law theorem numbered Proposition~1.9.  It neither assumes nor
proves first-passage stabilisation.  The separate edge combining this result
with Proposition~1.14 to obtain Proposition~1.11 must still be checked before
Theorem~1.6 or Theorem~1.3 is admitted as proved in this edition.  No
natural-density statement, convergence-to-one statement, or fixed universal
orbit bound follows from this proposition.
\end{nonclaim}
